We now shed light on the observations of the previous empirical results through a theoretical analysis. Our main result shows that DPT (under a slight modification to pretraining) essentially performs in-context posterior sampling (PS). PS is a generalization of Thompson Sampling for RL in MDPs. It maintains and samples from a posterior over tasks $\tau$ given historical data $D$ and executes optimal policies $\pistar_\tau$ (see Appendix~\ref{app:theory} for a formal outline). It is provably sample-efficient with online Bayesian regret guarantees~\citep{osband2013more}, but maintaining posteriors is generally computationally intractable. The ability for \macroName{} to perform PS in-context suggests a path towards computation- and provably sample-efficient RL with priors learned from the data.

\subsection{History-Dependent Pretraining and Assumptions}

We start with a modification to the pretraining of \macroName{}. Rather than conditioning only on $\squery$ and $D$ to predict $a^\star \sim \pistar_\tau(\cdot | \squery)$, we propose also conditioning on a sequence $\xi_{h} = (s_{1:h}, a_{1:h}^\star)$ where $s_{1:h} \sim \Sfrak_h \in \Delta(\Scal^h)$ is a distribution over sets of states, independent of $\tau$, and $a^\star_{h'} \sim \pistar_\tau(\cdot | s_{h'})$ for $h ' \in [h]$. Thus, we use $\pistar_\tau$ to label both the query state (which is the prediction label) and the sequence of states sampled from $\Sfrak_h$. Note that this does not require any environment interactions and hence no sampling from either $T_\tau$ or $R_\tau$.
At test-time at step $h$, this will allow us to condition on the history $\xi_{h - 1}$ of states that $M_{\theta}$ visits and the actions that it takes in those states. Formally, the learned $M_{\theta}$ is deployed as follows, given $D$.
% \begin{itemize}[leftmargin=1em,noitemsep]
    % \item 
    (1) At $h = 0$, initialize $\xi_0 = ()$ to be empty.
    % \item 
    (2) At step $h$, visit $s_h$ and find $a_h$ by sampling from  $M_{\theta}(\cdot | \squery, D, \xi_{h- 1})$.
    % \item 
    (3) Append $(s_h, a_h)$ to $\xi_{h -1}$ to get $\xi_h$.
% \end{itemize}
Note for bandits and contextual bandits ($H = 1$), there is no difference between this and the original pretraining procedure of prior sections because $\xi_0$ is empty. For MDPs, the original DPT can be viewed as a convenient approximation.

We now make several assumptions to simplify the analysis. First, assume $\Dquery$, $\pretraind$, and $\Sfrak$ have sufficient support such that all conditional probabilities of $P_{pre}$ are well defined. Similar to other studies of in-context learning~\citep{xie2021explanation}, we assume $M_{\theta}$ fits the pretraining distribution exactly with enough coverage and data, so that the focus of the analysis is just the in-context learning abilities.
\begin{assumption}\label{asmp:equality}
(Learned model is consistent). Let $M_{\theta}$ denote the pretrained model. For all $(\squery, D, \xi_h)$,
we have $P_{pre}(a | \squery, D, \xi_h) = M_{\theta}( a | \squery, D, \xi_h)$ for all $a \in \Acal$.
\end{assumption}
To provide some cursory justification, if $M_{\theta}$ is the global minimizer of~\eqref{eq:pretrain-obj}, then $\E_{P_{pre}} \| P_{pre}(\cdot |\squery,  D, \xi_h)  - M_{\theta}(\cdot |  \squery, D, \xi_h) \|_1^2 \rightarrow  0$ as the number of pretraining samples $N \rightarrow \infty$ with high probability for transformer model classes of bounded complexity (see Proposition~\ref{prop:mle}). Approximate versions of the above  %support
 assumptions are easily possible but obfuscate the key elements of the analysis.
We also assume that the in-context dataset $D \sim \pretraind$ is compliant~\citep{jin2021pessimism}, meaning that the actions from $D$ can depend only on the observed history and not additional confounders. Note that this still allows $\pretraind$ to be very general --- it could be generated randomly or from adaptive algorithms like PPO or TS.

\begin{definition}[Compliance]\label{def:compliant}
    The in-context dataset distribution $\pretraind(\cdot; \tau)$ is \emph{compliant} if, for all $i \in [n]$, the $i$th action of the dataset, $a_i$, is conditionally independent of $\tau$ given the $i$th state $s_i$ and partial dataset, $D_{i - 1}$, so far. In other words, the distribution $\pretraind(a_i | s_i, D_{i - 1}; \tau)$ is invariant to $\tau$.
\end{definition}

Generally, $\pretraind$ can influence $M_{\theta}$. In Proposition~\ref{prop:invariance}, we show that all compliant $\pretraind$ form a sort of equivalence class that generate the same $M_{\theta}$. For the remainder, we assume all $\pretraind$ are compliant.

\subsection{Main Results}

\paragraph{Equivalence of DPT and PS.} We now state our main result which shows that the trajectories generated by a pretrained $M_{\theta}$ will follow the same distribution as those from a well-specified PS algorithm. In particular, let PS use the well-specified prior $\pretraint$. Let $\tau_c$ be an arbitrary task. 
Let $P_{ps}(\cdot \ |\  D, \tau_c)$ and $P_{M_{\theta}}(\cdot \ | \ D, \tau_c)$ denote the distributions over trajectories $\xi_H \in (\Scal \times \Acal)^{H}$ generated from running PS and $M_{\theta}(\cdot | \cdot, D, \cdot)$, respectively, in task $\tau_c$ given historical data $D$.

\begin{restatable}[DPT $\iff$ PS]{theorem}{thmmain}
% \begin{theorem}[DPT $\iff$ PS]
\label{thm:main}
Let the above assumptions hold. Then, $P_{ps}(\xi_H  \ | \ D, \tau_c) = P_{M_{\theta}} (\xi_H \ | \ D, \tau_c)$ for all trajectories $\xi_H$.
% \end{theorem}
\end{restatable}


\paragraph{Regret implications.} To see this result in action, let us specialize to the finite MDP setting~\citep{osband2013more}. Suppose we pretrain $M_{\theta}$ on a distribution $\pretraint$ over MDPs with $S := |\Scal|$ and $A := |\Acal|$. Let $\pretraind$ be constructed by uniform sampling $(s_i, a_i)$ and observing $(r_i, s_i')$ for $i \in [KH]$. Let $\E\left[r_h | s_h, a_h\right] \in [0, 1]$. And let $\Dquery$ and $\Sfrak_h$ be uniform over $\Scal$ and $\Scal^h$ (for all $h$) respectively. Finally, let $\testt$ be the distribution over test tasks with the same cardinalities.
For a task $\tau$, define the online cumulative regret of DPT over $K$ episodes as   $\regret_\tau(M_{\theta}) := \sum_{k \in [K]} V_\tau(\pistar_\tau) - V_\tau(\hat \pi_k)$ where $\hat \pi_k(\cdot |s_h) = M_{\theta} (\cdot  | s_h, D_{(k - 1)}, \xi_{h - 1})$ and $D_{(k)}$ contains the first $k$ episodes collected from $\hat \pi_{1:k}$.
\begin{restatable}[Finite MDPs]{corollary}{cormdp}
% \begin{corollary}[Finite MDPs]
\label{cor:mdp}
    Suppose that $\sup_{\tau} {\testt(\tau)}/{\pretraint(\tau)} \leq \Ccal$ for some $\Ccal > 0$. For the above MDP setting, the pretrained model $M_{\theta}$ satisfies
   $\E_{\testt} \left[ \regret_\tau(M_{\theta}) \right]  \leq \widetilde{\Ocal} (  \Ccal H^{3/2}S\sqrt{ A K }) $.
% \end{corollary}
\end{restatable}

A similar analysis due to \cite{russo2014learning} allows us to prove why pretraining on (latently) linear bandits can lead to substantial empirical gains, even when the in-context datasets are generated by algorithms unaware of this structure. We observed this empirically in Section~\ref {sec:bandits}. Consider a similar setup as there where $\Scal$ is a singleton, $\Acal$ is finite but large, $\theta_\tau \in \RR^d$ is sampled as $\theta_{\tau} \sim \Ncal(0, I/d)$,   $\phi: \Acal \to \RR^d$ is a fixed feature map with $\sup_{a \in \Acal} \| \phi(a) \|_2 \leq 1$, and the reward of $a \in \Acal$ in task $\tau$ is distributed as $\Ncal(\dotp{\theta_\tau}{\phi(a)}, 1)$.  This time, we let $\pretraind(\cdot ; \tau)$ be given by running Thompson Sampling with Gaussian priors and likelihood functions on $\tau$.

\begin{restatable}[Latent representation learning in linear bandits]{corollary}{corlin}
% \begin{corollary}[Latent representation learning in linear bandits]
\label{cor:lin} For $\testt = \pretraint$ in the above linear bandit setting, $M_{\theta}$ satisfies
   $\E_{\testt} \left[ \regret_\tau(M_{\theta}) \right]  \leq \widetilde{\Ocal} ( d\sqrt{   K }) $.
% \end{corollary}
\end{restatable}


This significantly improves over the $\widetilde{\Ocal}(\sqrt{|\Acal| K } )$ upper regret bound for TS that does not leverage the linear structure. This highlights how DPT can have provably tighter upper bounds on future bandit problems than the algorithms used to generate its (pretraining) data. Note that if there is additional structure in the tasks which yields a tighter regret bound (for example if there are only a small finite number of known MDPs in the possible distribution), that may further improve performance, such as by removing the dependence on the problem finite state, action or full d-dimensional representation. % **EB: someone should cehck this 

\paragraph{Invariance of $M_{\theta}$ to compliant $\pretraind$.} Our final result sheds light on how $\pretraind$ impacts the final DPT behavior $M_{\theta}$. Combined with Assumption~\ref{asmp:equality}, $M_{\theta}$ is invariant to $\pretraind$ satisfying Definition~\ref{def:compliant}.
\begin{restatable}{proposition}{propinvariance}
% \begin{proposition}
\label{prop:invariance}
Let $P_{pre}^1$ and $P_{pre}^2$ be pretraining distributions that differ only by their in-context dataset distributions, denoted by $\pretraind^1$ and $\pretraind^2$. If $\pretraind^1$ and $\pretraind^2$ are compliant with the same support, then $P_{pre}^1(a^\star| \squery, D, \xi_h) = P_{pre}^2(a^\star| \squery, D, \xi_h)$ for all $a^\star, \squery, D, \xi_h$.
% \end{proposition}
\end{restatable}


That is, if we generate in-context datasets $D$ by running various algorithms that depend only on the observed data in the current task, we will end up with the same $M_{\theta}$. For example, TS could be used for $\pretraind^1$ and PPO for $\pretraind^2$. Expert-biased datasets discussed in Section~\ref{sec:bandits} violate Definition~\ref{def:compliant}, since privileged knowledge of $\tau$ is being used. This helps explain our empirical results that pretraining on expert-biased datasets leads to a qualitatively different learned model at test-time.